
et le
composé de ces homomorphismes pour \(0\leq i\leq n-1\) définit un homomorphisme (7.17.4.1), dont on peut montrer (procédant de façon analogue qu'en 7.3) que c'est le même que celui défini précédemment au moyen d'un éclatement. L'avantage de la deuxième méthode est cependant qu'on a même ici des homomorphismes de la forme (7.16.1) :
\[
  K_\bullet(X_{\bar t_{i+1}})\longrightarrow K_\bullet(X_{\bar t_i}),
\]
(avec des choix cohérents convenables des clôtures algébriques \(k(\bar t_i)\) des \(k(t_i)\)), qui permettent de définir par composition un homomorphisme
\begin{equation}\tag{7.17.4.2}
  K_\bullet(X_{\bar t})\longrightarrow K_\bullet(X_{\bar s}),
\end{equation}
qui dépend sans doute de façon essentielle du choix des \(S_i\), c'est-à-dire des \(t_i\).

\textbf{7.17.5.} On peut généraliser la construction d'un homomorphisme de spécialisation (7.17.4.2) au cas où on ne suppose plus que \(S\) soit régulier en \(s\) et \(t\) point maximal, pourvu qu'on dispose de la résolution des singularités pour le sous-schéma fermé intègre \(T\) de \(S\) de point générique \(t\), ou simplement pour son localisé en \(s\)~; cette hypothèse est vérifiée en particulier si \(\underline{O}_{S,s}\) est excellent de caractéristique nulle, grâce à Hironaka. En effet, on peut alors trouver un morphisme \(S'\to \Spec(\underline{O}_{T,s})\), avec \(S'\) régulier, et prenant un point fermé de la fibre \(S'_s\), on peut appliquer ce qui précède à \(X'=X\times_S S'\) sur \(S'\), le point générique \(t'\) de \(S'\) et son point fermé \(s'\). D'où un homomorphisme de la forme (7.17.4.2), puisque \(k(s')\) est une extension finie de \(k(s)\) et a donc même clôture algébrique.

\textbf{7.17.6.} Dans toutes les constructions envisagées dans la section 7.17, on s'est ramené pour la construction d'un homomorphisme de spécialisation au
